% Source: https://artofproblemsolving.com/community/c3h3162362_2017_lmt_spring_guts_round_p2536__lexington_mathematical_tournament
% Stable identity: aops-topic-3162362-q29
The intercept equation is $\dfrac{20}{a}+\dfrac{17}{b}=1 $. Cauchy--Schwarz gives $(a+b)(\dfrac{20}{a}+\dfrac{17}{b})\ge(\sqrt{20}+\sqrt{17})^2=37+4\sqrt{85}$. Equality holds for $ a=20+2\sqrt{85}$ and $ b=17+2\sqrt{85}$, which satisfy the intercept equation. Hence $\boxed{37+4\sqrt{85}}$.
\begin{asy}
usepackage("amsmath");
size(7cm);
real a=20+2sqrt(85),b=17+2sqrt(85);draw((0,0)--(a+3,0),Arrow);draw((0,0)--(0,b+3),Arrow);draw((a,0)--(0,b),blue);dot("$(20,17)$",(20,17),NE);label("$(a,0)$",(a,0),S);label("$(0,b)$",(0,b),W);
\end{asy}
